\section{未来、Scaling Limits 与 Continual Learning}

最后一章从应用愿景转向能力缺口。课程没有把更大模型视为唯一答案，而是依次追问可控性与记忆、具身交互与效率、可解释性、规模扩张的饱和点，以及部署后的持续适应。这里最重要的教学动作是给“持续学习”严格边界：RAG、长上下文和定期 fine-tuning 都可能更新行为，却不等于模型在部署中永久获得新能力。

判断一种方案是否真的完成持续学习，可以连续问四个问题：新知识写进了参数还是外部存储，更新是否跨会话保留，旧任务能力是否经过回归测试，以及错误更新能否定位并回滚。缺少其中任何一项，系统也许仍然有用，但更适合称为检索增强、临时上下文适应或离线再训练。这个区分会直接改变评测设计：不能只展示一次回答变好，还要测持久性、遗忘、冲突知识与更新成本。

\subsection{未来应用与缺失能力}

前面的应用案例已经展示 Transformer 能跨文本、图像和脑信号迁移；结尾的 future slides 把这种通用性外推到更长视频、医疗、环境和交互世界。读这类愿景页时要把“技术上可构造 demo”“在 benchmark 上有效”和“可承担真实责任”分成三层。每向后一层，都需要新的数据、系统验证、监管和人类责任边界。

Generalist agents、视频生成、数字人、诊断、环境监测和游戏 NPC 都需要 Transformer 与外部系统结合。应用清单支持的是机会空间，不证明模型已满足可靠性、实时性、隐私和责任要求。


\lecturefigure{slide-111.jpg}{未来应用：generalist agents、视频、数字人、医疗、环境与交互世界}{CS25 V5 official deck, p.~111}


\lecturefigure{slide-112.jpg}{缺失能力 I：计算效率、controllability、brain alignment 与跨域泛化}{CS25 V5 official deck, p.~112}


\lecturefigure{slide-113.jpg}{缺失能力 II：外部记忆、自我改进、自治、情感与价值对齐}{CS25 V5 official deck, p.~113}

\begin{knowledgebox}{深读：把 future wishlist 改写成可验收能力}
``Infinite memory'' 应改写为在长时间跨度内检索正确经历、处理冲突并遵守删除请求；``autonomy'' 应改写为在权限和预算内完成长任务、遇到不确定性及时升级；``emotional intelligence'' 需要在跨文化场景中识别情绪、避免操纵并校准信心；``embodiment'' 需要感知—行动闭环、物理安全和 sim-to-real；``self-improvement'' 则需要可测的新能力、受控遗忘和 rollback。把宏大名词拆成 observable criteria，团队才能设计 benchmark、红队与上线门槛。Slides 111--113 支持研究议程，不支持对 AGI 时间线或产品 readiness 的直接结论。
\end{knowledgebox}

\teachervoice{00:48:03--00:50:01，课堂先列应用，再列 missing ingredients：memory、continual learning、embodiment、autonomy 与 ethics。讲义提醒：愿景和能力缺口必须相邻书写，避免 roadmap 变成营销清单。}

\subsection{Minified models 与 interpretability}

On-device LLM 追求更低 memory、latency 和 energy。参数存储近似
\[
M_{\text{weights}}=N_{\text{params}}\cdot b/8,
\]
其中 $b$ 是 bits per parameter；实际还要加 KV cache、activation 与 runtime buffer。Quantization、distillation、sparsity 和 architecture redesign 各自改变质量与硬件友好性。


\lecturefigure{slide-114.jpg}{Minified/on-device LLM：低延迟、隐私、成本与可访问性}{CS25 V5 official deck, p.~114}


\lecturefigure{slide-115.jpg}{Interpretability：内部机制、failure analysis 与透明度挑战}{CS25 V5 official deck, p.~115}

\begin{knowledgebox}{深读：Interpretability 的目标需要先写清楚}
Mechanistic interpretability 试图定位内部 circuit；feature attribution 解释某次输出受哪些输入影响；behavioral testing 通过反事实 prompt 描述系统规律；monitoring 则在部署中发现异常。四者回答不同问题。一个 probe 能预测某层是否包含概念，不代表模型实际使用该概念；一个 saliency map 看起来合理，也不保证 causal faithfulness；一个 neuron edit 改掉事实，可能破坏无关知识。若目标是安全控制，还要问解释是否可重复、是否能提前预测 failure、能否指导 mitigation，以及 adversary 是否可绕过。Slide 115 的价值是把“参数太多、难以枚举”变成研究议程；它不支持把任一可视化工具当作透明度完成证明。
\end{knowledgebox}

\begin{warningbox}{小模型不自动更安全，可解释也不等于可控}
参数少可能便于部署和实验，却仍能幻觉、泄露或被提示攻击。Interpretability 方法若只在 toy tasks 上稳定，也不能直接作为 production safety guarantee；需要与 red teaming、permissions、monitoring 和 human oversight 联合。
\end{warningbox}

\subsection{Scaling limits：边际收益、数据与遗忘}

课程提出 pretraining scaling 可能出现 diminishing returns，并指出 post-training 也受 base model knowledge、reward signal 和 catastrophic forgetting 限制。若能力 $P(C)$ 随 compute $C$ 增长，可用局部斜率
\[
\alpha(C)=\frac{d\log P}{d\log C}
\]
监控边际收益；$\alpha$ 下降不等于 scaling “终结”，但说明数据、objective、architecture 或 evaluation 可能成为新瓶颈。


\lecturefigure{slide-116.jpg}{Limits of Scaling：数据、成本、知识饱和与边际收益}{CS25 V5 official deck, p.~116}


\lecturefigure{slide-117.jpg}{What is next：新架构、数据、硬件、理论与 adaptability}{CS25 V5 official deck, p.~117}

\begin{knowledgebox}{读图：Scaling limit 不是一条单变量曲线}
Slides 116--117 同时列出 data scarcity、hardware cost、architecture、post-training 和理论理解。读图时先区分\textbf{capability saturation}、\textbf{benchmark saturation} 与\textbf{economic saturation}：分数不涨可能因为评测封顶，成本不可接受也可能在能力仍提升时出现。局部 slope 下降支持寻找新瓶颈，却不能证明所有规模轴都失效。更好的实验是固定一种资源，单独扫描 model、data、compute 与 inference budget。
\end{knowledgebox}

\teachervoice{00:50:05--00:54:40，讲者把 on-device、interpretability、pretraining saturation、post-training limits 与 catastrophic forgetting 连在一起。课堂观点是“下一步不只更大”，但它仍是 2025 年的研究判断，不是已证明的 scaling 终点。}

\subsection{Continual learning 的严格定义}

Continual/lifelong learning 指模型在部署后从新经验持续更新能力，同时保持旧能力。至少要同时测量
\[
\text{Plasticity}=P_{\text{new}}^{\text{after}}-P_{\text{new}}^{\text{before}},\qquad
\text{Forgetting}=P_{\text{old}}^{\text{before}}-P_{\text{old}}^{\text{after}}.
\]
只提高新任务而严重遗忘旧任务，不是成功；只把新事实放进向量库，也没有更新参数能力。


\lecturefigure{slide-118.jpg}{Continual learning：部署后从反馈、经验与新数据持续改进}{CS25 V5 official deck, p.~118}


\lecturefigure{slide-119.jpg}{真正 lifelong learning 需要哪些机制：梯度、记忆、模块与 meta-learning}{CS25 V5 official deck, p.~119}

\teachervoice{00:54:43--00:57:10，老师明确排除 RAG、in-context learning、一次 fine-tuning、周期 retraining 和单纯 distillation，把 continual learning 定义为部署后永久、基础性的能力更新。这个定义较严格，讲义保留为 instructor position。}

\subsection{Model editing 与多种近似路线}

ROME 通过 causal tracing 定位与事实预测相关的层/表示，再做 rank-one update；MEMIT 扩展到多事实批量编辑。它们适合知识更新，却不等于学会数学或规划技能，也面临 locality、generalization 和 related-fact consistency。


\lecturefigure{slide-120.jpg}{Model editing：ROME 的 causal localization 与 rank-one update}{CS25 V5 official deck, p.~120}

其他路线包括 mistake-driven updates、lifelong mixture-of-experts、compressed prompt memory 和 progressive prompts。它们分别修改权重、增加 expert、维护外部 prompt memory 或组合 soft prompts；应按“更新位置”而不是名称分类。


\lecturefigure{slide-121.jpg}{Continual-learning works I：ROME/MEMIT、mistake updates 与 lifelong MoE}{CS25 V5 official deck, p.~121}


\lecturefigure{slide-122.jpg}{Continual-learning works II：prompt memory、CLOB 与 progressive prompts}{CS25 V5 official deck, p.~122}

\begin{knowledgebox}{术语消化：持续学习方法按“更新在哪里”分类}
RAG 更新外部知识库；CLOB 与 progressive prompts 更新提示记忆；ROME 和 MEMIT 修改局部权重；持续微调（continual fine-tuning）更新全局权重；lifelong MoE 增加或冻结 expert；meta-learning 则更新适应规则。Slides 120--122 的名称很多，但用“更新发生在哪里”这一坐标即可消化。只有当新经验能永久改变目标能力、同时旧能力保持并可审计回滚时，才接近严格 continual learning；否则更准确的说法可能只是 retrieval、adaptation 或 memory augmentation。
\end{knowledgebox}

\teachervoice{00:57:10--01:00:16，课堂逐项区分 ROME/MEMIT、CEM、lifelong MoE、CLOB 与 progressive prompts，并给出个人判断：真正 continual learning 可能需要更新模型“brain/weights”。讲义同时保留外部记忆路线的工程价值。}

\begin{table}[H]
\centering
\small
\begin{tabularx}{\textwidth}{L{0.18\textwidth} L{0.20\textwidth} X X}
\toprule
机制 & 更新位置 & 能解决什么 & 主要缺口 \\
\midrule
RAG/prompt memory & 外部上下文 & 新知识访问、可追溯 & 不改变基础能力 \\
Model editing & 局部权重/表示 & 事实修正 & 相关事实、技能与稳定性 \\
Continual fine-tuning & 全局权重 & 新域与技能 & 遗忘、成本、数据治理 \\
Lifelong MoE & 新/旧 experts & 模块化扩展 & 路由、容量与专家漂移 \\
Progressive prompts & soft prompts & 任务适配 & 组合爆炸、非基础更新 \\
Meta-learning & 更新规则 & 更快适应 & 训练复杂、稳定性难 \\
\bottomrule
\end{tabularx}
\caption{持续适应方法按更新位置和能力边界分类。}
\end{table}

\begin{lstlisting}[caption={Continual-learning evaluation 的最小协议}]
for task in task_stream:
    before = evaluate(model, old_tasks + [task])
    model = update(model, task.data, bounded_compute=True)
    after = evaluate(model, old_tasks + [task])
    log(plasticity=after[task]-before[task])
    log(forgetting=before[old_tasks]-after[old_tasks])
    audit(edit_locality, calibration, safety, rollback)
\end{lstlisting}

\subsection{本章小结}

未来路线同时要求更小、更透明、更能适应，而不是只扩大参数。Scaling 的局部边际收益可能下降；continual learning 则要求在获得新能力时控制遗忘。RAG、model editing、MoE expansion 与 weight updates 各自更新不同状态，必须用 plasticity、forgetting、locality 和系统成本共同评估。


\begin{importantbox}{跨章节证据等级：不要把所有 slide 写成同一种事实}
本讲至少包含四类证据。第一类是\textbf{定义与计算机制}，如 Q/K/V、DPO loss、ViT patch 和 cross-attention，可由公式与实现检查；第二类是\textbf{特定实验结果}，如 TinyDialogues、two-phase pretraining 和 fMRI downstream performance，只在给定模型、数据和 protocol 下成立；第三类是\textbf{工程启发式}，如先用小模型筛 blend、用 verifier 锚定 self-critique，需要在新环境复测；第四类是\textbf{研究判断与愿景}，如 scaling saturation、真正 continual learning 需要改权重、未来 generalist agents，应标明 speaker opinion 与日期。读者在引用任何结论前，都应先给它归类，再决定需要复现代码、补充实验、系统压测还是仅作为研究问题。这样才能让 overview 成为决策地图，而不是把定义、结果和预测混成一份权威清单。
\end{importantbox}

